% Apunts de teoria en PDF · 3r ESO · Unitat 11: Cossos geomètrics
% Mateix estil que els apunts de Batxillerat (franja i caixes de teoria en blau, exemples en roig), però més visuals,
% curts i amb dibuixos, per a l'alumnat de 3r ESO. Es compila amb pdflatex dues vegades (per als enllaços de l'índex),
% amb logo-fentmates.pdf (_eines/logo/pdf/horitzontal-clar.pdf) al costat.
\documentclass[11pt,a4paper]{article}
\usepackage[utf8]{inputenc}
\usepackage[T1]{fontenc}
\usepackage[catalan]{babel}
\usepackage[sfdefault,book,lining]{FiraSans}
\usepackage[italic,noplusnominus,noequal]{mathastext}
\linespread{1.18}
\usepackage[a4paper,top=2.4cm,bottom=2.2cm,left=1.6cm,right=1.6cm,headheight=1.1cm,headsep=0.5cm,footskip=1cm]{geometry}
\usepackage{amsmath,amssymb,cancel}
\usepackage{icomma}
\usepackage{graphicx}
\usepackage[table]{xcolor}
\usepackage{tikz}\usetikzlibrary{calc,positioning,arrows.meta,decorations.pathreplacing,shapes.geometric}
\usepackage{fontawesome5}
\usepackage{ccicons}
\usepackage{fancyhdr}
\usepackage{enumitem}
\usepackage{array}
\usepackage{needspace}
\usepackage[most]{tcolorbox}
\usepackage{titlesec}
\usepackage[hidelinks]{hyperref}
\hypersetup{pdftitle={Apunts de teoria. Unitat 11: Cossos geomètrics. 3r ESO},pdfsubject={Matemàtiques 3r ESO},pdfkeywords={poliedres, fórmula d'Euler, prismes, piràmides, cilindre, con, esfera, coordenades geogràfiques, àrea, volum, Cavalieri}}

% colors de la web (bloc d'Àlgebra: blau)
\definecolor{fmnavy}{HTML}{1B2748}
\definecolor{fmblue}{HTML}{3D6AA0}
\definecolor{fmbluel}{HTML}{4E81BD}
\definecolor{fmyellow}{HTML}{FFC94D}
\definecolor{fmgold}{HTML}{E9A919}
\definecolor{fmaccent}{HTML}{FF7A50}
\definecolor{fmdef}{HTML}{1D4F73}\definecolor{fmdefbg}{HTML}{EEF4F9}
\definecolor{fmex}{HTML}{A3203A}\definecolor{fmexbg}{HTML}{FBF0F2}
\definecolor{fmsf}{HTML}{C2410C}
\definecolor{fmok}{HTML}{15803D}
\definecolor{fmmuted}{HTML}{5D6685}
\definecolor{fmsurface}{HTML}{EEF1F8}
\definecolor{fmA}{HTML}{2F6FD6}\definecolor{fmB}{HTML}{EA6A10}\definecolor{fmR}{HTML}{E11D48}\definecolor{fmK}{HTML}{16A34A}\definecolor{fmV}{HTML}{7C3AED}

% lletres
\newcommand{\comic}{\fontfamily{ComicNeue-TLF}\selectfont}
\newcommand{\heavy}{\fontseries{eb}\selectfont}
\newcommand{\semi}{\fontseries{sb}\selectfont}
\DeclareUnicodeCharacter{2212}{-}

\setlength{\parindent}{0pt}
\setlength{\parskip}{7pt}
\setlength{\lineskip}{4pt}\setlength{\lineskiplimit}{4pt}
\setlength{\emergencystretch}{3em}
\setlist{itemsep=3pt,topsep=3pt,leftmargin=1.4em}
\setcounter{secnumdepth}{1}
\renewcommand{\thesection}{11.\arabic{section}}
\titleformat{\section}[hang]{\color{fmnavy}\fontsize{23}{27}\heavy\selectfont}{\tcbox[on line,colback=fmblue,colframe=fmblue,arc=3mm,boxsep=2pt,left=6pt,right=6pt,top=3pt,bottom=3pt]{\color{white}\thesection}}{0.45em}{}[\vspace{5pt}{\color{fmblue}\titlerule[2pt]}]
\titlespacing*{\section}{0pt}{4pt}{6pt}
\titleformat{\subsection}{\color{fmblue}\Large\semi}{}{0pt}{\faCaretRight\ }
\titlespacing*{\subsection}{0pt}{8pt}{2pt}

% frase de cada apartat, «a mà»
\newcommand{\frase}[1]{\par\vspace{-2pt}{\comic\Large\color{fmbluel}\faCommentDots[regular]\ \ #1\par}\vspace{2pt}}
\newcommand{\seccio}[3]{\Needspace{16\baselineskip}\section{#2}\hypertarget{#1}{}\frase{#3}}
% notes a mà dins dels exemples
\newcommand{\nota}[1]{{\comic\small\color{fmex!85!black}#1}}
\newcommand{\notab}[1]{{\comic\small\color{fmblue}#1}}
% bé / malament
\newcommand{\be}[1]{{\color{fmok}\faCheckCircle}\ \ensuremath{#1}}
\newcommand{\mal}[1]{{\color{fmR}\faTimesCircle}\ {\color{fmR}\ensuremath{#1}}}
% pas numerat
\newcommand{\pas}[1]{\tikz[baseline=(p.base)]\node[circle,fill=fmblue,text=white,inner sep=1.2pt,minimum size=1.25em,font=\small\heavy](p){#1};}
% etiqueta de color
\newcommand{\chip}[2]{\tcbox[on line,colback=#1!14,colframe=#1!60,arc=2mm,boxsep=1pt,left=5pt,right=5pt,top=1pt,bottom=1pt]{#2}}
\newcommand{\fmkey}[1]{\colorbox{fmyellow!55}{\strut #1}}
\newcommand{\cA}[1]{{\color{fmA}#1}}\newcommand{\cB}[1]{{\color{fmB}#1}}\newcommand{\cK}[1]{{\color{fmK}#1}}\newcommand{\cR}[1]{{\color{fmR}#1}}\newcommand{\cV}[1]{{\color{fmV}#1}}

% caixes
\tcbset{fmbase/.style={enhanced,arc=3mm,boxrule=1.2pt,left=4mm,right=4mm,top=2mm,bottom=2.5mm,before skip=8pt,after skip=8pt,fonttitle=\heavy\large,before upper=\raggedright}}
\newtcolorbox{idea}[1][Idea clau]{fmbase,colback=fmdefbg,colframe=fmdef,coltitle=white,title={\faKey\ \ #1}}
\newtcolorbox{exemple}[1][Exemple]{fmbase,colback=fmexbg,colframe=fmex,colbacktitle=fmexbg,coltitle=fmex,title={\faPencil*\ \ #1},titlerule=0.5pt,titlerule style={fmex!40,dashed}}
\newtcolorbox{compte}[1][Compte!]{fmbase,colback=fmsf!6,colframe=fmsf,coltitle=white,fonttitle=\comic\bfseries\Large,title={\faExclamationTriangle\ \ #1}}
\newtcolorbox{truc}[1][Truc]{fmbase,colback=fmyellow!20,colframe=fmgold,colbacktitle=fmyellow,coltitle=fmnavy,fonttitle=\comic\bfseries\Large,title={\faLightbulb[regular]\ \ #1}}

% dibuix de la targeta de la unitat, en blanc
\newcommand{\fmcard}{\begin{tikzpicture}[x=0.062cm,y=-0.062cm,white,line width=3pt,line cap=round,line join=round]
  \begin{scope}[transparency group,opacity=0.15]
  \node[anchor=base west,inner sep=0pt,font=\fontsize{30}{30}\selectfont] at (14,44) {$C+V=A+2$};
  \draw (130,105) -- (170,105) -- (170,65) -- (130,65) -- cycle; \draw (130,65) -- (150,50) -- (190,50) -- (170,65); \draw (190,50) -- (190,90) -- (170,105);
  \draw (215,105) ellipse (18 and 6); \draw (197,105) -- (215,55) -- (233,105);
  \end{scope}
\end{tikzpicture}}
\newcommand{\fmbrand}{\raisebox{-0.25\height}{\includegraphics[height=0.5cm]{logo-fentmates.pdf}}}

% regla de Ruffini: #1 = a, #2 = coeficients, #3 = fila del mig (sense el primer), #4 = fila de baix, #5 = nombre de coeficients
\newcommand{\ruffini}[5]{\begin{tikzpicture}[x=1.25cm,y=0.85cm,every node/.style={font=\large}]
  \foreach \v [count=\i] in {#2} \node (t\i) at (\i,0) {$\v$};
  \foreach \v [count=\i] in {#3} {\pgfmathtruncatemacro\j{\i+1}\node (m\j) at (\j,-1) {$\v$};}
  \foreach \v [count=\i] in {#4} \node (b\i) at (\i,-2.1) {$\v$};
  \node[text=fmB] at (0,-1) {$#1$};
  \draw[line width=1.1pt,fmnavy] (0.45,0.5)--(0.45,-1.5) (0.45,-1.5)--(#5+0.5,-1.5);
  \node[draw=fmR,line width=1.3pt,rounded corners=2mm,minimum size=0.75cm] at (#5,-2.1) {};
  \draw[-{Latex[length=2.2mm]},fmK,line width=1.1pt] (t1.south) -- (b1.north);
  \pgfmathtruncatemacro\last{#5-1}
  \foreach \i in {1,...,\last} {\pgfmathtruncatemacro\j{\i+1}
    \draw[-{Latex[length=2.2mm]},fmB,line width=1.1pt,shorten >=1pt] (b\i.north east) -- (m\j.south west);
    \draw[-{Latex[length=2.2mm]},fmA!70,line width=1pt,dashed,shorten >=2pt] ([xshift=8pt]t\j.south) -- ([xshift=8pt]b\j.north);}
\end{tikzpicture}}

% capçalera i peu
\pagestyle{fancy}\fancyhf{}
\fancyhead[L]{\small\fmbrand}
\fancyhead[R]{\small\color{fmmuted}3r ESO · Unitat 11: Cossos geomètrics}
\fancyfoot[L]{\scriptsize\color{fmmuted}\href{https://creativecommons.org/licenses/by-nc-sa/4.0/deed.ca}{\ccbyncsa}\ \ FentMates · fentmates.org · CC BY-NC-SA 4.0}
\fancyfoot[R]{\small\color{fmmuted}\thepage}
\renewcommand{\headrulewidth}{0.4pt}\renewcommand{\headrule}{\color{fmmuted!40}\hrule height \headrulewidth}
\fancypagestyle{primera}{\fancyhf{}\renewcommand{\headrulewidth}{0pt}
  \fancyfoot[L]{\scriptsize\color{fmmuted}\href{https://creativecommons.org/licenses/by-nc-sa/4.0/deed.ca}{\ccbyncsa}\ \ FentMates · fentmates.org · CC BY-NC-SA 4.0}
  \fancyfoot[R]{\small\color{fmmuted}\thepage}}

% targeta de l'índex
\newcommand{\itile}[4]{\begin{tcolorbox}[enhanced,colback=white,colframe=fmblue!35,arc=3mm,boxrule=1pt,boxsep=0pt,left=6pt,right=6pt,top=5pt,bottom=5pt,before skip=0pt,after skip=7pt]
  \hyperlink{#1}{\begin{minipage}[c]{0.9cm}\centering{\color{fmblue}\LARGE #3}\end{minipage}\hspace{6pt}%
  \begin{minipage}[c]{0.78\linewidth}{\color{fmblue}\heavy #2}\\[-2pt]{\color{fmnavy}\semi\large #4}\end{minipage}}\end{tcolorbox}}

% taula amb vores arrodonides: el contingut es retalla amb la forma (sense @{} a les vores ni entre columnes)
\newsavebox\tabclipbox
\newcommand{\taulaclip}[3][3mm]{\sbox\tabclipbox{#3}\begin{tikzpicture}[baseline=(tcn.base)]
  \begin{scope}\clip[rounded corners=#1] (0,-\dp\tabclipbox) rectangle (\wd\tabclipbox,\ht\tabclipbox);
  \node[inner sep=0pt,outer sep=0pt,anchor=base west] (tcn) at (0,0) {\usebox\tabclipbox};\end{scope}
  \draw[#2,rounded corners=#1] (0,-\dp\tabclipbox) rectangle (\wd\tabclipbox,\ht\tabclipbox);
\end{tikzpicture}}
\newcommand{\taularod}[3][fmblue]{\taulaclip{#1,line width=1.2pt}{\renewcommand{\arraystretch}{1.5}\begin{tabular}{#2}#3\end{tabular}}}
% balança
\newcommand{\balanca}[3]{\begin{tikzpicture}[scale=#3,transform shape,every node/.style={anchor=center}]
  \fill[fmnavy] (-0.45,-1.9) rectangle (0.45,-1.75);
  \fill[fmnavy!80] (-0.12,-1.75) rectangle (0.12,0);
  \fill[fmgold] (0,0.05) circle (0.13);
  \draw[fmnavy,line width=2.4pt,line cap=round] (-2,0) -- (2,0);
  \foreach \s in {-1,1} {
    \draw[fmmuted,line width=0.8pt] (\s*2,0) -- (\s*2-0.75,-1.05) (\s*2,0) -- (\s*2+0.75,-1.05);
    \fill[fmyellow!70,draw=fmgold,line width=1pt] (\s*2-0.95,-1.05) arc[start angle=180,end angle=360,x radius=0.95,y radius=0.32] -- cycle;}
  \node[font=\Large,fill=fmA!15,draw=fmA,rounded corners=1.5mm,inner sep=3pt] at (-2,-0.68) {$#1$};
  \node[font=\Large,fill=fmB!15,draw=fmB,rounded corners=1.5mm,inner sep=3pt] at (2,-0.68) {$#2$};
\end{tikzpicture}}
% eixos amb quadrícula (dins d'un tikzpicture): xmin, xmax, ymin, ymax
\newcommand{\eixos}[4]{\draw[fmmuted!22,line width=0.4pt] (#1,#3) grid (#2,#4);
  \draw[-{Latex[length=2mm]},fmmuted,line width=0.8pt] (#1-0.2,0)--(#2+0.4,0) node[right,font=\small]{$x$};
  \draw[-{Latex[length=2mm]},fmmuted,line width=0.8pt] (0,#3-0.2)--(0,#4+0.4) node[above,font=\small]{$y$};}
% taula de traduccions (diccionari) amb les vores arrodonides
\newcommand{\dic}[2]{#1 & {\large$#2$}\\}
\newcommand{\dicratio}{0.62}
\newcommand{\dichead}{\begin{tabular}{@{}>{\raggedright}p{\dimexpr\dicratio\linewidth-12pt}>{\centering\arraybackslash}p{\dimexpr\linewidth-\dicratio\linewidth-12pt}@{}}\faComment[regular]\ Ho diem així\dots & \dots\ així\end{tabular}}
\newenvironment{dicbox}{\begin{tcolorbox}[enhanced,arc=3mm,boxrule=1.2pt,colframe=fmblue,colback=white,coltitle=white,fonttitle=\heavy,
  title={\dichead},left=0pt,right=0pt,top=0pt,bottom=1.5mm,boxsep=0pt,lefttitle=6pt,righttitle=6pt,before skip=0pt,after skip=0pt]\rowcolors{1}{white}{fmsurface}%
  \begin{tabular}{@{\hspace{6pt}}>{\raggedright}p{\dimexpr\dicratio\linewidth-12pt}>{\centering\arraybackslash}p{\dimexpr\linewidth-\dicratio\linewidth-12pt}@{\hspace{6pt}}}}
  {\end{tabular}\end{tcolorbox}}
% resum
\newcommand{\rtile}[3]{\begin{tcolorbox}[enhanced,colback=white,colframe=fmblue!45,arc=3mm,boxrule=1pt,left=3mm,right=3mm,top=2.5mm,bottom=2.5mm,fontupper=\large,
  title={\tcbox[on line,colback=fmblue,colframe=fmblue,arc=2mm,boxsep=0pt,left=4pt,right=4pt,top=1pt,bottom=1pt]{\color{white}\small\heavy #1}\ \ #2},coltitle=fmnavy,colbacktitle=fmdefbg,fonttitle=\semi,toptitle=0.5mm,bottomtitle=0.5mm]#3\end{tcolorbox}}
% portada: #1 subtítol a mà
\newcommand{\portada}[1]{\thispagestyle{primera}
\begin{tikzpicture}[remember picture,overlay]
  \fill[fmblue] (current page.north west) rectangle ([yshift=-9.4cm]current page.north east);
  \begin{scope}
    \clip (current page.north west) rectangle ([yshift=-9.4cm]current page.north east);
    \node[anchor=north west,inner sep=0pt] at ([xshift=6.6cm,yshift=-1.3cm]current page.north west) {\fmcard};
  \end{scope}
  \node[anchor=north west,inner sep=0pt] at ([xshift=1.6cm,yshift=-1.1cm]current page.north west) {\tcbox[colback=white,colframe=white,arc=2.5mm,boxsep=0pt,left=7pt,right=7pt,top=5pt,bottom=5pt]{\includegraphics[height=0.75cm]{logo-fentmates.pdf}}};
\end{tikzpicture}
\vspace*{0.55cm}
\tcbox[colback=white!15!fmblue,colframe=white!15!fmblue,arc=2.5mm,boxsep=0pt,left=6pt,right=6pt,top=3pt,bottom=3pt]{\color{white}\small\bfseries Apunts de teoria · 3r ESO · Matemàtiques}

\vspace{2pt}
{\raggedright\color{white}\fontsize{30}{36}\heavy\selectfont{\color{fmyellow}Unitat 11}\\[2pt]Cossos geomètrics\par}
\vspace{8pt}
{\color{white}\comic\Large #1\par}
\vspace*{\dimexpr 7.4cm-\pagetotal\relax}
}
\newcommand{\llegenda}{\par\vspace{16pt}
{\color{fmnavy}\heavy\Large\faPalette\ \ Com es llegeixen aquests apunts}\par\smallskip
\begin{tcbraster}[raster columns=4,raster equal height,raster column skip=8pt,raster before skip=12pt]
\tcbox[colback=fmdefbg,colframe=fmdef,arc=3mm,boxrule=1.2pt,halign=center,width=\linewidth]{{\color{fmdef}\Large\faKey}\\[2pt]{\color{fmdef}\heavy Idea clau}\\[-1pt]{\small el que cal saber}}
\tcbox[colback=fmexbg,colframe=fmex,arc=3mm,boxrule=1.2pt,halign=center,width=\linewidth]{{\color{fmex}\Large\faPencil*}\\[2pt]{\color{fmex}\heavy Exemple}\\[-1pt]{\small pas a pas}}
\tcbox[colback=fmsf!6,colframe=fmsf,arc=3mm,boxrule=1.2pt,halign=center,width=\linewidth]{{\color{fmsf}\Large\faExclamationTriangle}\\[2pt]{\color{fmsf}\heavy Compte!}\\[-1pt]{\small errors típics}}
\tcbox[colback=fmyellow!20,colframe=fmgold,arc=3mm,boxrule=1.2pt,halign=center,width=\linewidth]{{\color{fmgold!80!black}\Large\faLightbulb[regular]}\\[2pt]{\color{fmnavy}\heavy Truc}\\[-1pt]{\small per a recordar-ho}}
\end{tcbraster}}
\newcommand{\resumtitol}{\clearpage
\hypertarget{resum}{}
{\color{fmnavy}\fontsize{23}{27}\heavy\selectfont\faClipboardList\ \ En resum}\par\vspace{5pt}{\color{fmblue}\titlerule[2pt]}
\frase{Tota la unitat en una pàgina.}}
% els 5 passos per a resoldre problemes: #2 = text del pas 2
\newcommand{\passosproblema}[2]{\begin{center}
\begin{tikzpicture}[s/.style={draw=fmblue,fill=fmdefbg,rounded corners=3mm,line width=1.2pt,minimum width=2.75cm,minimum height=2.45cm,align=center,font=\small},
  ic/.style={font=\LARGE,text=fmblue}]
  \foreach \i/\ic/\t in {1/\faBookOpen/{\textbf{Llegim}\\i entenem},2/\faQuestionCircle/{#1},3/\faEquals/{#2},4/\faCogs/{\textbf{Resolem}},5/\faCheckDouble/{\textbf{Comprovem}\\i responem}} {
    \node[s] (s\i) at (3.3*\i,0) {};
    \node[ic] at ([yshift=-0.5cm]s\i.north) {\ic};
    \node[align=center,font=\small,anchor=north] at ([yshift=-0.95cm]s\i.north) {\t};
    \node[circle,fill=fmblue,text=white,font=\small\heavy,inner sep=1.5pt,minimum size=1.3em] at (s\i.north west) {\i};}
  \foreach \i [evaluate=\i as \j using int(\i+1)] in {1,...,4} \draw[-{Latex[length=3mm]},fmmuted,line width=1.5pt] (s\i) -- (s\j);
\end{tikzpicture}
\end{center}}

\newcommand{\lo}{\mathopen{]}}\newcommand{\ro}{\mathclose{[}}
% punts ple i buit, nombres dels eixos i taula de valors
\newcommand{\pple}[2][fmblue]{\fill[#1] #2 circle (3.6pt);}
\newcommand{\nums}[4]{\foreach \k in {#1,...,#2} {\ifnum\k=0\else\node[below,font=\scriptsize,text=fmmuted] at (\k,0) {$\k$};\fi}
  \foreach \k in {#3,...,#4} {\ifnum\k=0\else\node[left,font=\scriptsize,text=fmmuted] at (0,\k) {$\k$};\fi}}
\newcommand{\tv}[3][fmblue]{\taulaclip{#1,line width=1.2pt}{\renewcommand{\arraystretch}{1.45}%
  \begin{tabular}{>{\columncolor{#1}\color{white}\heavy}c*{12}{w{c}{0.8cm}}}#2\\\rowcolor{#1!8}\cellcolor{#1}#3\end{tabular}}}
% etiqueta d'una recta o corba
\newcommand{\etq}[3]{\node[fill=white,inner sep=1.5pt,text=#1,font=\small] at #2 {#3};}

% ortoedre en perspectiva: #1 ample, #2 alt, #3 fondària (en perspectiva), #4 color
\newcommand{\ortoedre}[4]{\filldraw[fill=#4!18,draw=#4,line width=1.2pt] (0,0) rectangle (#1,#2);
  \filldraw[fill=#4!30,draw=#4,line width=1.2pt] (0,#2) -- ++(0.6*#3,0.4*#3) -- ++(#1,0) -- (#1,#2) -- cycle;
  \filldraw[fill=#4!45,draw=#4,line width=1.2pt] (#1,0) -- ++(0.6*#3,0.4*#3) -- ++(0,#2) -- (#1,#2) -- cycle;
  \draw[#4,dashed,line width=0.7pt] (0,0) -- ++(0.6*#3,0.4*#3) -- ++(#1,0) (0.6*#3,0.4*#3) -- ++(0,#2);}
% cilindre: #1 radi, #2 altura, #3 color
\newcommand{\cilindre}[3]{\fill[#3!18] (-#1,0) -- (-#1,#2) arc (180:360:#1 and 0.3*#1) -- (#1,0) arc (0:-180:#1 and 0.3*#1) -- cycle;
  \draw[#3,line width=1.2pt] (-#1,0) -- (-#1,#2) (#1,0) -- (#1,#2);
  \filldraw[fill=#3!35,draw=#3,line width=1.2pt] (0,#2) ellipse (#1 and 0.3*#1);
  \draw[#3,line width=1.2pt] (-#1,0) arc (180:360:#1 and 0.3*#1); \draw[#3,dashed,line width=0.7pt] (#1,0) arc (0:180:#1 and 0.3*#1);}
% con: #1 radi, #2 altura, #3 color
\newcommand{\con}[3]{\fill[#3!18] (-#1,0) -- (0,#2) -- (#1,0) arc (0:-180:#1 and 0.3*#1) -- cycle;
  \draw[#3,line width=1.2pt] (-#1,0) -- (0,#2) -- (#1,0); \draw[#3,line width=1.2pt] (-#1,0) arc (180:360:#1 and 0.3*#1); \draw[#3,dashed,line width=0.7pt] (#1,0) arc (0:180:#1 and 0.3*#1);}

\begin{document}
\portada{Prismes, piràmides, cilindres, cons i esferes: superfícies i volums.}

{\color{fmnavy}\heavy\Large\faMapSigns\ \ Què aprendràs?}\par\smallskip
\begin{minipage}[t]{0.485\linewidth}
\itile{sec1}{11.1}{\faCube}{Poliedres. Fórmula d'Euler}
\itile{sec2}{11.2}{\faBox}{Prismes: elements i àrea}
\itile{sec3}{11.3}{\faIcon{caret-up}}{Piràmides i troncs}
\itile{sec4}{11.4}{\faIcon{database}}{Cilindre i con}
\end{minipage}\hfill
\begin{minipage}[t]{0.485\linewidth}
\itile{sec5}{11.5}{\faGlobeEurope}{L'esfera. La Terra}
\itile{sec6}{11.6}{\faCubes}{Volum de prismes i cilindres}
\itile{sec7}{11.7}{\faIcon{ice-cream}}{Volum de piràmides, cons i esferes}
\itile{sec8}{11.8}{\faPuzzlePiece}{Cossos compostos}
\itile{resum}{\phantom{11.8}}{\faClipboardList}{En resum}
\end{minipage}
\llegenda

% ============================================================ 11.1
\clearpage
\seccio{sec1}{Poliedres. Fórmula d'Euler}{Cossos limitats per polígons.}

\vfill
\begin{minipage}[c]{0.4\linewidth}\centering
\begin{tikzpicture}[scale=0.85]
  \ortoedre{3}{2.2}{2}{fmA}
  \node[font=\small\heavy,text=fmA] at (1.5,1.1) {cara};
  \draw[fmR,line width=2.4pt] (3,0) -- (3,2.2); \node[right,font=\small\heavy,text=fmR] at (4.25,0.3) {aresta};
  \draw[-{Latex[length=2mm]},fmR,line width=0.8pt] (4.25,0.3) -- (3.1,0.6);
  \fill[fmK] (4.2,3) circle (3.5pt); \node[right,font=\small\heavy,text=fmK!80!black] at (4.3,3.1) {vèrtex};
\end{tikzpicture}
\end{minipage}\hfill
\begin{minipage}[c]{0.57\linewidth}
\begin{idea}[Poliedre]
Un \textbf{poliedre} és un cos limitat per \textbf{polígons}:\\[3pt]
\cA{\textbf{cares}}: els polígons; \ \cR{\textbf{arestes}}: els costats de les cares; \ \cK{\textbf{vèrtexs}}: on es troben les arestes.\\[4pt]
És \textbf{convex} si totes les cares es poden recolzar sobre una taula; si no, \textbf{còncau}.
\end{idea}
\end{minipage}

\vfill
\begin{tcolorbox}[enhanced,colback=fmyellow!18,colframe=fmgold,arc=4mm,boxrule=1.6pt,before upper=\centering,top=2mm,bottom=2mm]
{\Large Fórmula d'Euler: \quad $C+V=A+2$}\par\smallskip
{\small en tot poliedre convex \quad prisma pentagonal: \ $7+10=15+2$ \ \be{}}
\end{tcolorbox}

\vfill
\subsection{Els cinc poliedres regulars}
\begin{minipage}[c]{0.62\linewidth}\centering\small
\renewcommand{\arraystretch}{1.3}
\taularod{>{\heavy}l<{\hspace{4pt}}>{\hspace{4pt}}c<{\hspace{4pt}}>{\hspace{4pt}}c<{\hspace{4pt}}>{\hspace{4pt}}c<{\hspace{4pt}}>{\hspace{4pt}}l<{\hspace{4pt}}>{\hspace{4pt}}c<{\hspace{4pt}}}{
\rowcolor{fmblue}\color{white}Poliedre & \color{white}\heavy $C$ & \color{white}\heavy $V$ & \color{white}\heavy $A$ & \color{white}\heavy Cares & \color{white}\heavy Per vèrtex\\
Tetraedre & 4 & 4 & 6 & triangles & 3\\
\rowcolor{fmsurface}Cub & 6 & 8 & 12 & quadrats & 3\\
Octaedre & 8 & 6 & 12 & triangles & 4\\
\rowcolor{fmsurface}Dodecaedre & 12 & 20 & 30 & pentàgons & 3\\
Icosaedre & 20 & 12 & 30 & triangles & 5
}
\end{minipage}\hfill
\begin{minipage}[c]{0.35\linewidth}
\begin{truc}[Per què només cinc?]
En cada vèrtex, els angles de les cares han de sumar \textbf{menys de $360^\circ$}. Tres hexàgons ja fan $360^\circ$: queda pla.
\end{truc}
\end{minipage}

\vfill
\begin{exemple}[Exemples]
\begin{minipage}[t]{0.48\linewidth}
Piràmide hexagonal: \ $C=7$, \ $V=7$, \ $A=12$\\ $7+7=12+2$ \ \be{}\\[4pt]
9 cares i 16 arestes: \ $9+V=18\ \to\ \boxed{V=9}$
\end{minipage}\hfill
\begin{minipage}[t]{0.49\linewidth}
Prisma de 30 arestes: \ $3n=30\to n=10$\\
base \textbf{decàgon}: \ $C=\boxed{12}$, \ $V=\boxed{20}$
\end{minipage}
\end{exemple}

\vfill
% ============================================================ 11.2
\clearpage
\seccio{sec2}{Prismes: elements i àrea}{Dues bases iguals i paral·leles, i cares laterals rectangulars.}

\vfill
\begin{minipage}[c]{0.36\linewidth}\centering
\begin{tikzpicture}[scale=0.75]
  \ortoedre{3}{2.4}{1.6}{fmB}
  \draw[fmR,line width=1.5pt] (0,0) -- (3.96,3.04);
  \draw[fmK,line width=1.3pt,dashed] (0,0) -- (3.96,0.64);
  \node[font=\small\heavy,text=fmR] at (1.6,1.7) {$D$};
  \node[below,font=\small] at (1.5,0) {$a$}; \node[right,font=\small] at (3.5,0.3) {$b$}; \node[right,font=\small] at (3.96,1.9) {$c$};
\end{tikzpicture}\\
{\comic\small\color{fmmuted}ortoedre i la seua diagonal}
\end{minipage}\hfill
\begin{minipage}[c]{0.6\linewidth}
\begin{idea}[Prisma]
Dues \textbf{bases} iguals i paral·leles, i \textbf{cares laterals} que són paral·lelograms. L'\textbf{altura} $h$ és la distància entre les bases.\\[3pt]
\textbf{Recte}: cares laterals rectangulars. \ \textbf{Regular}: recte i amb bases regulars.\\[3pt]
\textbf{Ortoedre}: una caixa. \ \textbf{Cub}: totes les cares quadrades.\\[4pt]
Diagonal de l'ortoedre: \ $\boxed{D=\sqrt{a^2+b^2+c^2}}$
\end{idea}
\end{minipage}

\vfill
\begin{minipage}[c]{0.44\linewidth}\centering
\begin{tikzpicture}[scale=0.42]
  \foreach \i in {0,...,5} \filldraw[fill=fmA!20,draw=fmA,line width=1pt] (\i*1.3,0) rectangle ++(1.3,4);
  \filldraw[fill=fmB!30,draw=fmB,line width=1pt] (2.6,4) -- ++(1.3,0) -- ++(0.65,1.126) -- ++(-0.65,1.126) -- ++(-1.3,0) -- ++(-0.65,-1.126) -- cycle;
  \filldraw[fill=fmB!30,draw=fmB,line width=1pt] (2.6,0) -- ++(1.3,0) -- ++(0.65,-1.126) -- ++(-0.65,-1.126) -- ++(-1.3,0) -- ++(-0.65,1.126) -- cycle;
  \draw[decorate,decoration={brace,amplitude=4pt,mirror},line width=0.8pt] (0,-0.2) -- (2.4,-0.2) node[midway,below=4pt,font=\scriptsize] {};
  \node[right,font=\small] at (7.8,2) {$h$};
\end{tikzpicture}\\
{\comic\small\color{fmmuted}desenvolupament d'un prisma hexagonal}
\end{minipage}\hfill
\begin{minipage}[c]{0.53\linewidth}
\begin{tcolorbox}[enhanced,colback=fmyellow!18,colframe=fmgold,arc=4mm,boxrule=1.6pt,before upper=\centering,top=2mm,bottom=2mm]
{\large$A_L=P_{\text{base}}\cdot h$}\par\medskip
{\large$A_T=A_L+2\,A_{\text{base}}$}\par\smallskip
{\small les cares laterals fan un rectangle de base el perímetre}
\end{tcolorbox}
\end{minipage}

\vfill
\begin{exemple}[Exemple: \ prisma hexagonal regular de $l=4$ cm i $h=10$ cm]
Apotema de la base: \ $a=\sqrt{4^2-2^2}=\sqrt{12}\approx3{,}46$ \qquad $A_L=24\cdot10=240$ cm$^2$\\[3pt]
$A_{\text{base}}=\dfrac{24\cdot3{,}46}{2}\approx41{,}57$ \qquad $A_T\approx240+2\cdot41{,}57=\boxed{323{,}14\text{ cm}^2}$
\end{exemple}

\vfill
\begin{exemple}[Exemples]
\begin{minipage}[t]{0.48\linewidth}
Ortoedre $8\times5\times3$:\\ $2\,(40+24+15)=\boxed{158\text{ cm}^2}$\\[4pt]
Cub de 6 cm: \ $6\cdot36=\boxed{216\text{ cm}^2}$
\end{minipage}\hfill
\begin{minipage}[t]{0.49\linewidth}
Pintar parets i sostre d'una habitació $5\times4\times2{,}5$ m (porta $0{,}8\times2$, finestra $1{,}5\times1{,}2$):\\
$45+20-1{,}6-1{,}8=61{,}6$ m$^2$; \ pots de 12 m$^2$: \boxed{6\text{ pots}}
\end{minipage}
\end{exemple}

\vfill
% ============================================================ 11.3
\clearpage
\seccio{sec3}{Piràmides i troncs de piràmide}{Cares laterals triangulars que es troben en el vèrtex.}

\vfill
\begin{minipage}[c]{0.4\linewidth}\centering
\begin{tikzpicture}[scale=1.05]
  \coordinate (A) at (0,0); \coordinate (B) at (3,0); \coordinate (C) at (4,0.8); \coordinate (D) at (1,0.8);
  \coordinate (O) at (2,0.4); \coordinate (V) at (2,3.4); \coordinate (M) at (2.5,0);
  \fill[fmA!15] (A) -- (B) -- (V) -- cycle; \fill[fmA!28] (B) -- (C) -- (V) -- cycle;
  \draw[fmA,line width=1.2pt] (A) -- (B) -- (C) (A) -- (V) -- (B) (C) -- (V);
  \draw[fmA,dashed,line width=0.7pt] (A) -- (D) -- (C) (D) -- (V);
  \fill[fmK!25] (O) -- (V) -- (1.5,0) -- cycle;
  \draw[fmR,line width=1.5pt] (O) -- (V) node[midway,left,font=\small\heavy] {$h$};
  \draw[fmB,line width=1.5pt] (O) -- (1.5,0) node[midway,below left=-2pt,font=\small\heavy] {$a_b$};
  \draw[fmK,line width=1.5pt] (1.5,0) -- (V) node[midway,left=1pt,font=\small\heavy] {$a_p$};
\end{tikzpicture}
\end{minipage}\hfill
\begin{minipage}[c]{0.57\linewidth}
\begin{idea}[Piràmide regular]
Base regular i cares laterals triangles isòsceles iguals.\\[3pt]
\cK{\textbf{Apotema de la piràmide} $a_p$}: l'altura de cada cara lateral.\\
\cB{\textbf{Apotema de la base} $a_b$}. \ No les confongues!\\[4pt]
Formen un triangle rectangle amb l'altura:\\[2pt]
{\centering\Large$a_p^2=h^2+a_b^2$\par}
\end{idea}
\end{minipage}

\vfill
\begin{tcolorbox}[enhanced,colback=fmyellow!18,colframe=fmgold,arc=4mm,boxrule=1.6pt,before upper=\centering,top=2mm,bottom=2mm]
{\Large$A_L=\dfrac{P_{\text{base}}\cdot a_p}{2}\qquad\qquad A_T=A_L+A_{\text{base}}$}\par\smallskip
{\small el desenvolupament és la base i $n$ triangles de base $l$ i altura $a_p$}
\end{tcolorbox}

\vfill
\begin{exemple}[Exemple: \ piràmide quadrangular regular de $l=6$ cm i $h=4$ cm]
$a_b=3$, \ $a_p=\sqrt{4^2+3^2}=5$ cm \qquad $A_L=\dfrac{24\cdot5}{2}=60$ cm$^2$ \qquad $A_T=60+36=\boxed{96\text{ cm}^2}$
\end{exemple}

\vfill
\begin{minipage}[c]{0.34\linewidth}\centering
\begin{tikzpicture}[scale=0.7]
  \filldraw[fill=fmB!20,draw=fmB,line width=1.2pt] (0,0) -- (4,0) -- (3,2) -- (1,2) -- cycle;
  \draw[fmK,line width=1.4pt] (1,2) -- (1,0) node[midway,left,font=\small\heavy] {$h$};
  \draw[fmR,line width=1.4pt] (0,0) -- (1,2);
  \node[below,font=\scriptsize] at (0.5,0) {$\frac{l-l'}{2}$};
  \node[above,font=\small,text=fmR] at (0.3,1.1) {$a_p$};
\end{tikzpicture}
\end{minipage}\hfill
\begin{minipage}[c]{0.63\linewidth}
\begin{tcolorbox}[fmbase,colback=fmdefbg,colframe=fmdef,coltitle=white,title={Tronc de piràmide}]
Tallem una piràmide amb un pla paral·lel a la base: dues bases semblants i cares laterals \textbf{trapezis}.\\[3pt]
{\centering$A_L=\dfrac{(P+P')\cdot a_p}{2}\qquad A_T=A_L+A_B+A_{B'}$\par}
\end{tcolorbox}
\end{minipage}

\vfill
\begin{exemple}[Exemple: \ tronc quadrangular regular, arestes de les bases 10 cm i 4 cm, altura 4 cm]
$a_p=\sqrt{4^2+3^2}=5$ \qquad $A_L=\dfrac{(40+16)\cdot5}{2}=140$ \qquad $A_T=140+100+16=\boxed{256\text{ cm}^2}$
\end{exemple}

\vfill
% ============================================================ 11.4
\clearpage
\seccio{sec4}{Cilindre i con}{Cossos que ixen de fer girar una figura plana.}

\vfill
\begin{minipage}[c]{0.5\linewidth}\centering
\begin{tikzpicture}[scale=0.75]
  \begin{scope}\cilindre{1}{2.8}{fmA}\draw[fmR,line width=1.4pt] (0,2.8) -- (1,2.8) node[midway,above,font=\small\heavy] {$r$}; \draw[fmK,line width=1.4pt,dashed] (0,0) -- (0,2.8); \node[left,font=\small\heavy,text=fmK!80!black] at (0,1.4) {$h$};\end{scope}
  \begin{scope}[xshift=3.6cm]\con{1.1}{2.8}{fmB}\fill[fmK!25] (0,0) -- (1.1,0) -- (0,2.8) -- cycle;\draw[fmK,line width=1.4pt] (0,0) -- (0,2.8) node[midway,left,font=\small\heavy] {$h$};\draw[fmR,line width=1.4pt] (0,0) -- (1.1,0) node[midway,below,font=\small\heavy] {$r$};\draw[fmV,line width=1.6pt] (1.1,0) -- (0,2.8) node[midway,right,font=\small\heavy] {$g$};\end{scope}
\end{tikzpicture}
\end{minipage}\hfill
\begin{minipage}[c]{0.47\linewidth}
\begin{idea}[Cossos de revolució]
Un rectangle que gira sobre un costat fa un \textbf{cilindre}; un triangle rectangle que gira sobre un catet, un \textbf{con}.\\[4pt]
En el con, la \cV{\textbf{generatriu} $g$} va del vèrtex a la vora de la base:\\[2pt]
{\centering\Large$g^2=h^2+r^2$\par}
\end{idea}
\end{minipage}

\vfill
\begin{minipage}[c]{0.48\linewidth}\centering
\begin{tikzpicture}[scale=0.38]
  \filldraw[fill=fmA!20,draw=fmA,line width=1pt] (0,0) rectangle (6.28,3);
  \filldraw[fill=fmA!35,draw=fmA,line width=1pt] (3.14,4) circle (1); \filldraw[fill=fmA!35,draw=fmA,line width=1pt] (3.14,-1) circle (1);
  \node[below,font=\scriptsize] at (5,0) {$2\pi r$}; \node[right,font=\scriptsize] at (6.28,1.5) {$h$};
  \begin{scope}[xshift=9.5cm,yshift=0.5cm]
  \filldraw[fill=fmB!20,draw=fmB,line width=1pt] (0,0) -- (-30:4) arc (-30:60:4) -- cycle;
  \filldraw[fill=fmB!35,draw=fmB,line width=1pt] (15:5) circle (1);
  \node[font=\scriptsize] at (20:2) {$g$};
  \end{scope}
\end{tikzpicture}\\
{\comic\small\color{fmmuted}desenvolupaments: rectangle i dos cercles; sector i cercle}
\end{minipage}\hfill
\begin{minipage}[c]{0.49\linewidth}\centering
\renewcommand{\arraystretch}{1.6}
\taularod{>{\heavy}l<{\hspace{5pt}}>{\hspace{5pt}}c<{\hspace{5pt}}>{\hspace{5pt}}c<{\hspace{5pt}}}{
\rowcolor{fmblue} & \color{white}\heavy Àrea lateral & \color{white}\heavy Àrea total\\
Cilindre & $2\pi rh$ & $2\pi rh+2\pi r^2$\\
\rowcolor{fmsurface}Con & $\pi rg$ & $\pi rg+\pi r^2$
}
\end{minipage}

\vfill
\begin{exemple}[Exemple: \ con de $r=3$ cm i $h=4$ cm]
$g=\sqrt{16+9}=5$ \qquad $A_L=\pi\cdot3\cdot5=15\pi\approx47{,}12$ cm$^2$ \qquad $A_T=15\pi+9\pi=24\pi\approx\boxed{75{,}40\text{ cm}^2}$
\end{exemple}

\vfill
\begin{exemple}[Exemples]
\begin{minipage}[t]{0.48\linewidth}
Cilindre de $r=5$, $h=12$:\\ $A_T=120\pi+50\pi=170\pi\approx\boxed{534{,}07\text{ cm}^2}$\\[4pt]
Con de 8 cm de diàmetre i $h=10$:\\ $g=\sqrt{100+16}\approx10{,}77$, \ $A_T\approx\boxed{185{,}6\text{ cm}^2}$
\end{minipage}\hfill
\begin{minipage}[t]{0.49\linewidth}
Llanda de 7,4 cm de diàmetre i 11 cm d'altura:\\
etiqueta $=A_L=2\pi\cdot3{,}7\cdot11\approx\boxed{255{,}7\text{ cm}^2}$\\
xapa $=A_T\approx255{,}7+2\pi\cdot3{,}7^2\approx\boxed{341{,}7\text{ cm}^2}$
\end{minipage}
\end{exemple}

\vfill
\begin{compte}[Generatriu, no altura]
L'àrea lateral del con usa la \textbf{generatriu} $g$. Si et donen l'altura, calcula primer $g$ amb Pitàgores.
\end{compte}

\vfill
% ============================================================ 11.5
\clearpage
\seccio{sec5}{L'esfera. La Terra}{Tots els punts a la mateixa distància del centre.}

\vfill
\begin{minipage}[c]{0.36\linewidth}\centering
\begin{tikzpicture}[scale=0.85]
  \shade[ball color=fmA!35] (0,0) circle (1.6);
  \draw[fmA,line width=1.3pt] (0,0) circle (1.6);
  \draw[fmB,line width=1.2pt] (-1.6,0) arc (180:360:1.6 and 0.45); \draw[fmB,dashed,line width=0.8pt] (1.6,0) arc (0:180:1.6 and 0.45);
  \draw[fmR,line width=1.3pt] (0,0) -- (1.6,0) node[midway,above,font=\small\heavy] {$r$}; \fill[fmnavy] (0,0) circle (1.5pt);
  \node[font=\scriptsize,text=fmB] at (0,-0.75) {cercle màxim};
\end{tikzpicture}
\end{minipage}\hfill
\begin{minipage}[c]{0.6\linewidth}
\begin{idea}[L'esfera]
La \textbf{superfície esfèrica} són els punts a distància $r$ del centre; l'\textbf{esfera} és el cos que tanca. Un pla pel centre la talla en un \textbf{cercle màxim}.\\[4pt]
No es pot desplegar sobre un pla (els mapes del món deformen). La seua àrea són \textbf{4 cercles màxims}:\\[2pt]
{\centering\fontsize{20}{24}\selectfont$A=4\pi r^2$\par}
\end{idea}
\end{minipage}

\vfill
\begin{truc}[Arquimedes]
L'àrea de l'esfera és igual a l'àrea lateral del cilindre que l'envolta: \ $2\pi r\cdot2r=4\pi r^2$. \ Va demanar que gravaren aquesta figura a la seua tomba.
\end{truc}

\vfill
\begin{minipage}[c]{0.36\linewidth}\centering
\begin{tikzpicture}[scale=0.85]
  \shade[ball color=fmK!30] (0,0) circle (1.6);
  \draw[fmnavy,line width=1.2pt] (0,0) circle (1.6);
  \foreach \l in {-50,-25,25,50} \draw[fmB,line width=0.8pt] ({-1.6*cos(\l)},{1.6*sin(\l)}) arc (180:360:{1.6*cos(\l)} and {0.4*cos(\l)});
  \draw[fmR,line width=1.4pt] (-1.6,0) arc (180:360:1.6 and 0.4); \node[right,font=\scriptsize,text=fmR] at (1.6,0) {equador};
  \foreach \a in {0.5,1.1} \draw[fmV,line width=0.8pt] (0,1.6) arc (90:270:{\a} and 1.6);
  \draw[fmV,line width=0.8pt] (0,1.6) -- (0,-1.6);
  \node[above,font=\scriptsize] at (0,1.6) {N}; \node[below,font=\scriptsize] at (0,-1.6) {S};
\end{tikzpicture}
\end{minipage}\hfill
\begin{minipage}[c]{0.6\linewidth}
\begin{tcolorbox}[fmbase,colback=fmdefbg,colframe=fmdef,coltitle=white,title={\faGlobeEurope\ \ La Terra ($r\approx6371$ km)}]
\cV{\textbf{Meridians}}: de pol a pol. \ \cB{\textbf{Paral·lels}}: perpendiculars a l'eix; el més gran és l'\cR{\textbf{equador}}.\\[3pt]
\textbf{Latitud}: angle amb l'equador, $0^\circ$ a $90^\circ$ N o S.\\
\textbf{Longitud}: angle amb Greenwich, $0^\circ$ a $180^\circ$ E o O.\\[3pt]
València: $39{,}5^\circ$ N, $0{,}4^\circ$ O. \ Cada $15^\circ$ de longitud, \textbf{1 hora} de diferència.
\end{tcolorbox}
\end{minipage}

\vfill
\begin{exemple}[Exemples]
Esfera de $r=6$ cm: \ $A=4\pi\cdot36=144\pi\approx\boxed{452{,}39\text{ cm}^2}$\\[5pt]
Esfera de 1256,64 cm$^2$: \ $r^2=\dfrac{1256{,}64}{4\pi}=100\ \to\ \boxed{r=10\text{ cm}}$\\[5pt]
La Terra: \ $4\pi\cdot6371^2\approx\boxed{510\text{ milions de km}^2}$; \ si el 71\,\% és aigua, terra ferma: \ $0{,}29\cdot510\approx\boxed{148\text{ milions de km}^2}$
\end{exemple}

\vfill
% ============================================================ 11.6
\clearpage
\seccio{sec6}{Volum. Prismes i cilindres}{L'espai que ocupa un cos: àrea de la base per altura.}

\vfill
\begin{tcbraster}[raster columns=2,raster equal height,raster column skip=10pt]
\begin{tcolorbox}[fmbase,colback=fmdefbg,colframe=fmdef,coltitle=white,title={\faCube\ \ Unitats de volum: de 1000 en 1000}]
{\centering$1\text{ m}^3=1000\text{ dm}^3$\qquad $1\text{ dm}^3=1000\text{ cm}^3$\par}
\end{tcolorbox}
\begin{tcolorbox}[fmbase,colback=fmK!8,colframe=fmK,coltitle=white,title={\faTint\ \ Capacitat}]
{\centering$1\text{ dm}^3=\boxed{1\text{ L}}$\qquad $1\text{ cm}^3=1\text{ mL}$\qquad $1\text{ m}^3=1000\text{ L}$\par}
\end{tcolorbox}
\end{tcbraster}

\vfill
\begin{minipage}[c]{0.42\linewidth}\centering
\begin{tikzpicture}[scale=0.55]
  \foreach \i in {0,...,6} \filldraw[fill=fmgold!50,draw=fmgold!80!black,line width=0.8pt] (0,\i*0.45) rectangle (2.4,\i*0.45+0.4);
  \foreach \i in {0,...,6} \filldraw[fill=fmgold!50,draw=fmgold!80!black,line width=0.8pt] ({4+0.35*\i},\i*0.45) rectangle ++(2.4,0.4);
  \node[below,font=\scriptsize] at (1.2,0) {recta}; \node[below,font=\scriptsize] at (5.2,0) {torta};
\end{tikzpicture}\\
{\comic\small\color{fmmuted}mateix volum}
\end{minipage}\hfill
\begin{minipage}[c]{0.55\linewidth}
\begin{idea}[Principi de Cavalieri]
Si dos cossos tenen la mateixa altura i, a qualsevol altura, les seccions tenen la \textbf{mateixa àrea}, tenen el \textbf{mateix volum}.
\end{idea}
\end{minipage}

\vfill
\begin{tcolorbox}[enhanced,colback=fmyellow!18,colframe=fmgold,arc=4mm,boxrule=1.6pt,before upper=\centering,top=2mm,bottom=2mm]
{\Large$V_{\text{prisma}}=A_{\text{base}}\cdot h\qquad\qquad V_{\text{cilindre}}=\pi r^2h$}\par\smallskip
{\small un ortoedre $a\times b\times c$ conté $a\cdot b\cdot c$ cubets; per Cavalieri, val per a tots els prismes}
\end{tcolorbox}

\vfill
\begin{exemple}[Exemples]
\begin{minipage}[t]{0.48\linewidth}
3,5 m$^3$ $=\boxed{3500\text{ L}}$ \quad 750 cm$^3=\boxed{0{,}75\text{ L}}$\\[4pt]
Cub de 7 cm: \ $7^3=\boxed{343\text{ cm}^3}$\\[4pt]
Cilindre $r=4$, $h=9$: \ $144\pi\approx\boxed{452{,}39\text{ cm}^3}$
\end{minipage}\hfill
\begin{minipage}[t]{0.49\linewidth}
Prisma hexagonal regular, $l=4$, $h=10$:\\ $A_{\text{base}}\approx41{,}57$ \ $\to$ \ $V\approx\boxed{415{,}7\text{ cm}^3}$\\[4pt]
Cilindre $r=3$, $h=10$: \ $90\pi\approx282{,}7$ cm$^3\approx0{,}28$ L
\end{minipage}
\end{exemple}

\vfill
\begin{exemple}[Exemple: \ una piscina de 25 m $\times$ 12 m $\times$ 1,8 m s'omple a 500 L per minut]
$V=25\cdot12\cdot1{,}8=540$ m$^3=\boxed{540\,000\text{ L}}$ \qquad $540\,000:500=1080$ min $=\boxed{18\text{ hores}}$
\end{exemple}

\vfill
% ============================================================ 11.7
\clearpage
\seccio{sec7}{Volum de piràmides, cons i esferes}{Una piràmide és la tercera part del prisma.}

\vfill
\begin{minipage}[c]{0.42\linewidth}\centering
\begin{tikzpicture}[scale=0.62]
  \begin{scope}\cilindre{1}{2.6}{fmA}\con{1}{2.6}{fmB}\end{scope}
  \node[font=\Large] at (2.1,1.3) {$=$};
  \begin{scope}[xshift=3.4cm]\con{0.55}{2.6}{fmB}\end{scope}
  \begin{scope}[xshift=4.6cm]\con{0.55}{2.6}{fmB}\end{scope}
  \begin{scope}[xshift=5.8cm]\con{0.55}{2.6}{fmB}\end{scope}
  \node[below,font=\scriptsize] at (4.6,-0.4) {tres cons plens};
\end{tikzpicture}
\end{minipage}\hfill
\begin{minipage}[c]{0.55\linewidth}
\begin{idea}[Un terç]
Amb una piràmide plena d'aigua, cal abocar \textbf{tres} vegades per a omplir el prisma de la mateixa base i altura. Igual amb el con i el cilindre.
\end{idea}
\end{minipage}

\vfill
\begin{tcolorbox}[enhanced,colback=fmyellow!18,colframe=fmgold,arc=4mm,boxrule=1.6pt,before upper=\centering,top=3mm,bottom=3mm]
{\Large$V_{\text{piràmide}}=\dfrac{A_{\text{base}}\cdot h}{3}\qquad V_{\text{con}}=\dfrac{\pi r^2h}{3}\qquad V_{\text{esfera}}=\dfrac43\pi r^3$}
\end{tcolorbox}

\vfill
\begin{truc}[Arquimedes, una altra vegada]
L'esfera és $\frac23$ del cilindre que l'envolta: \ $\frac23\cdot\pi r^2\cdot2r=\frac43\pi r^3$.
\end{truc}

\vfill
\begin{exemple}[Exemples]
Piràmide quadrangular de $l=6$ cm i $h=4$ cm: \ $V=\dfrac{36\cdot4}{3}=\boxed{48\text{ cm}^3}$\\[5pt]
Con de $r=5$ cm i $h=12$ cm: \ $V=\dfrac{\pi\cdot25\cdot12}{3}=100\pi\approx\boxed{314{,}16\text{ cm}^3}$\\[5pt]
Esfera de $r=3$ cm: \ $V=\frac43\pi\cdot27=36\pi\approx\boxed{113{,}10\text{ cm}^3}$ \qquad Pilota de $r=11$ cm: \ $\frac43\pi\cdot1331\approx\boxed{5{,}6\text{ L}}$
\end{exemple}

\vfill
\begin{exemple}[Exemple: \ con de generatriu 13 cm i radi 5 cm. Quina esfera té el mateix volum?]
$h=\sqrt{13^2-5^2}=12$ \ $\to$ \ $V=\dfrac{\pi\cdot25\cdot12}{3}=100\pi$ \qquad $\dfrac43\pi r^3=100\pi\ \to\ r^3=75\ \to\ r\approx\boxed{4{,}22\text{ cm}}$
\end{exemple}

\vfill
\begin{compte}[No t'oblides el 3]
Piràmide i con: \textbf{entre 3}. Si no el poses, calcules el volum del prisma o del cilindre.
\end{compte}

\vfill
% ============================================================ 11.8
\clearpage
\seccio{sec8}{Cossos compostos i problemes}{Descomponem en cossos coneguts i sumem (o restem).}

\vfill
\begin{minipage}[c]{0.3\linewidth}\centering
\begin{tikzpicture}[scale=0.62]
  \cilindre{1.2}{3.2}{fmA}
  \begin{scope}[yshift=3.2cm]\fill[fmB!25] (-1.2,0) -- (0,1.6) -- (1.2,0) -- cycle;\draw[fmB,line width=1.2pt] (-1.2,0) -- (0,1.6) -- (1.2,0);\end{scope}
  \node[right,font=\scriptsize] at (1.2,1.6) {6 m}; \node[right,font=\scriptsize] at (0.7,4) {3 m}; \node[below,font=\scriptsize] at (0.6,-0.4) {$r=2$};
\end{tikzpicture}\\
{\comic\small\color{fmmuted}sitja}
\end{minipage}\hfill
\begin{minipage}[c]{0.67\linewidth}
\begin{idea}[Cossos compostos]
\textbf{Volum}: sumem els volums dels cossos simples (i \textbf{restem} si hi ha forats).\\[3pt]
\textbf{Àrea}: sumem només les superfícies que queden \textbf{per fora}.
\end{idea}
\begin{exemple}[Exemple: \ sitja: cilindre + con]
$V=\pi\cdot4\cdot6+\dfrac{\pi\cdot4\cdot3}{3}=24\pi+4\pi=28\pi\approx\boxed{87{,}96\text{ m}^3}$
\end{exemple}
\end{minipage}

\vfill
\begin{exemple}[Exemple: \ un gelat: con de $r=3$ cm i $h=10$ cm, ple, amb mitja esfera al damunt]
$V=\dfrac{\pi\cdot9\cdot10}{3}+\dfrac12\cdot\dfrac43\pi\cdot27=30\pi+18\pi=48\pi\approx\boxed{150{,}8\text{ cm}^3}$
\end{exemple}

\vfill
\begin{exemple}[Exemple: \ un cub de fusta de 10 cm amb un forat cilíndric de 2 cm de radi que el travessa]
$V=1000-\pi\cdot4\cdot10=1000-40\pi\approx\boxed{874{,}3\text{ cm}^3}$\\[3pt]
Àrea: \ $600-2\cdot4\pi+2\pi\cdot2\cdot10=600+32\pi\approx\boxed{700{,}5\text{ cm}^2}$ \ \nota{6 cares, menys els 2 cercles, més l'interior}
\end{exemple}

\vfill
\begin{exemple}[Exemple: \ cubell en forma de tronc de con: fons $r=10$ cm, boca $r=15$ cm, altura 20 cm]
\begin{minipage}[c]{0.65\linewidth}
És un con gran menys un con menut. Per Tales: \ $\dfrac{x}{10}=\dfrac{x+20}{15}\ \to\ x=40$\\[3pt]
$V=\dfrac{\pi\cdot225\cdot60}{3}-\dfrac{\pi\cdot100\cdot40}{3}=\dfrac{9500\pi}{3}\approx9948$ cm$^3\approx\boxed{9{,}95\text{ L}}$
\end{minipage}\hfill
\begin{minipage}[c]{0.32\linewidth}\centering
\begin{tikzpicture}[scale=0.07]
  \fill[fmB!20] (-15,20) -- (-10,0) -- (10,0) -- (15,20) -- cycle;
  \draw[fmB,line width=1.2pt] (-15,20) -- (-10,0) (10,0) -- (15,20);
  \draw[fmB,line width=1.2pt] (0,20) ellipse (15 and 3); \draw[fmB,line width=1.2pt] (-10,0) arc (180:360:10 and 2);
  \draw[fmmuted,dashed] (-10,0) -- (0,-40) -- (10,0);
\end{tikzpicture}
\end{minipage}
\end{exemple}

\vfill
% ============================================================ resum
\resumtitol
\begin{tcbraster}[raster columns=2,raster equal height=rows,raster column skip=14pt,raster row skip=30pt,raster before skip=16pt]
\rtile{11.1}{Poliedres}{$C+V=A+2$.\\[1ex] Cinc regulars: tetraedre, cub, octaedre, dodecaedre, icosaedre.}
\rtile{11.2}{Prismes}{$A_L=P\cdot h$, \ $A_T=A_L+2A_B$.\\[1ex] Ortoedre: $D=\sqrt{a^2+b^2+c^2}$.}
\rtile{11.3}{Piràmides}{$A_L=\frac{P\cdot a_p}{2}$, \ $a_p^2=h^2+a_b^2$.\\[1ex] Tronc: cares trapezis.}
\rtile{11.4}{Cilindre i con}{Cilindre $A_L=2\pi rh$. Con $A_L=\pi rg$.\\[1ex] $g^2=h^2+r^2$.}
\rtile{11.5}{Esfera}{$A=4\pi r^2$.\\[1ex] Latitud (N/S) i longitud (E/O).}
\rtile{11.6}{Volum}{1 dm$^3=1$ L. \ $V=A_B\cdot h$, \ cilindre $\pi r^2h$.}
\rtile{11.7}{Un terç}{Piràmide $\frac{A_B\cdot h}{3}$, con $\frac{\pi r^2h}{3}$.\\[1ex] Esfera $\frac43\pi r^3$.}
\rtile{11.8}{Compostos}{Volums: sumem o restem.\\[1ex] Àrea: només per fora.}
\end{tcbraster}
\end{document}
